\begin{lemma}[Existing finite type of the actual original section algebra]
\label{scheme:existing-hironaka-original-section-finite-type}
\lean{AlgebraicGeometry.Scheme.Hom.finiteType_sectionsAlgebra}
\leanok
\uses{scheme:existing-hironaka-original-section-action}
Let $R$ be ANY commutative ring, $X$ ANY original scheme, and let
$f:X\to\operatorname{Spec}R$ be an actual locally finite type morphism.
For every actual affine open $U\subset X$, the SAME original ring
$\Gamma(X,U)$ is a finite type $R$-algebra for the already defined
$f.\mathrm{sectionsAlgebra}(U)$. Its scalar map is the original map
$R\xrightarrow{\Gamma\mathrm{SpecIso}^{-1}}
\Gamma(\operatorname{Spec}R,\top)\xrightarrow{f.\mathrm{appLE}}\Gamma(X,U)$.
There is no field, reducedness, integrality, smoothness, nonempty-open,
Noetherian or quasi-compactness assumption.
\end{lemma}
\begin{proof}
This is the existing Hironaka theorem for the existing named algebra
structure. The actual morphism locally of finite type gives finite type
of $f.\mathrm{appLE}(\top,U)$ between the original affine section rings.
The inverse spectrum global-section ring map is surjective because it
is an isomorphism, hence is finite type. Compose these two original
finite type ring maps. This composite is exactly the original scalar
map defining $f.\mathrm{sectionsAlgebra}(U)$, so it proves the stated
finite type algebra property. This records the original proof and
introduces no new Lean theorem or mathematical producer.
\end{proof}
